% !TEX root = documentation.tex


\titlehead{BSc Computational and Data Science\\CDS-309 Hyperautomation und RPA\\Dozenten: Markus Müntener, Matthias Müller\hfill}
\title{Automatisierung eines Moodle-basierten Lernplanungsprozesses mit UiPath und generativer KI}
\author[1,*]{Alessio De Icco}
\author[2,*]{Florin Gartmann}
\author[3,*]{Michal Karczmarzyk}
\author[4,*]{Oliver Schütz}
\affil[1]{Fachhochschule Graubünden}
\affil[*]{E-Mail Adressen: michal.karczmarzyk@stud.fhgr.ch, oliver.schuetz@stud.fhgr.ch, alessio.deicco@stud.fhgr.ch, florin.gartmann@stud.fhgr.ch}
\date{\today}
\maketitle

\begin{abstract}
Diese Arbeit untersucht die Automatisierung eines Moodle-basierten Lernplanungsprozesses mit UiPath, Moodle-Datenzugriff und generativer KI. Ausgangspunkt ist das Problem, dass Studierende Kursinformationen, Aufgaben, Deadlines und Semesterunterlagen häufig manuell in verschiedenen Moodle-Kursen suchen und anschliessend für die Erstellung eines Lernplans aufbereiten müssen.

Ziel des Moodle-to-LLM Study Planners ist es, diese wiederkehrenden Arbeitsschritte zu automatisieren. UiPath steuert den Prozess, liest relevante Moodle-Informationen aus, speichert sie strukturiert zwischen und übergibt sie an die Gemini API, die daraus einen Lernplan erstellt. Damit verbindet die Lösung klassische RPA mit API-Nutzung und GenAI zu einem Hyperautomation-Ansatz.

Die Analyse zeigt, dass der Prozess technisch gut automatisierbar ist, da er digital, wiederholbar und datengetrieben abläuft. Herausforderungen bestehen vor allem in uneinheitlichen Moodle-Strukturen, unterschiedlichen PDF-Formaten, API-Fehlern und der Validierung von LLM-Antworten. Durch strukturierte Datenverarbeitung, Logging und Retry-Mechanismen können diese Risiken reduziert werden. Insgesamt zeigt die Arbeit, dass der gewählte Prozess sowohl fachlich relevant als auch technisch und ökonomisch für eine RPA-Automatisierung geeignet ist.
\end{abstract}

\clearpage